\documentclass[10pt]{article}
\usepackage[utf8]{inputenc}
\usepackage[T1]{fontenc}
\usepackage[margin=1in, left=1.25in, right=1.25in, top=1in, bottom=1in]{geometry}

\pagestyle{empty}

\begin{document}
\begin{center}
{\Large Research highlights}
\end{center}

\begin{itemize}
  \setlength{\itemsep}{0.35em}
  \item Introduces a bi-level reinforcement learning (RL) framework for rolling bus timetable optimization and station holding control.
  \item Maps upper-level planned-headway adjustments to rolling planned dispatch headways and launch times.
  \item Couples asynchronous dispatch planning and holding via a dynamic timetable, holding feedback, and upper-only hindsight credit.
  \item Reports delay-aware service metrics, including holding time, onboard time, total passenger time, and dispatch lateness.
  \item Achieves lower mean composite cost than target-headway-only, fixed-timetable, classical holding, and same-simulator RL proxy baselines in the reported ten-seed comparison.
\end{itemize}

\end{document}
